%% ma: L10-B13-0048
%% lop: 10
%% chuong: 5
%% bai: 13
%% dang: 10.13.8
%% loai: DS
%% muc: [NB, TH, NB, TH]
%% dap_an: "SSSĐ"
%% nguon: "chuyen-de-cac-so-dac-trung-cua-mau-so-lieu-khong-ghep-nhom-toan-10-chuong-trinh-moi.pdf (GV. Phan Nhật Linh) (P02), Bài 02 Dạng 1 Phần II Câu 3"
%% kien_thuc: "Lớp 10 Bài 13 (số trung bình, trung vị)"
%% trang_thai: chinh-thuc
%% ghi_chu: "Đáp án gốc ý c) Đúng là sai (lời giải gốc sắp xếp nhầm, số lớn nhất là 67). Tự duyệt theo yêu cầu giáo viên 22/10/2026."
\begin{ex}
Một cửa hàng vật liệu xây dựng thống kê số bao xi măng bán ra trong $23$ ngày cuối năm. Kết quả như sau:
\begin{center}
\begin{tabular}{|c|c|c|c|c|c|c|c|c|c|c|c|}\hline
$47$ & $54$ & $43$ & $50$ & $61$ & $36$ & $65$ & $54$ & $50$ & $43$ & $62$ & $59$\\\hline
$36$ & $45$ & $45$ & $33$ & $53$ & $67$ & $21$ & $45$ & $50$ & $36$ & $58$ & \\\hline
\end{tabular}
\end{center}
Xét tính đúng sai của các khẳng định sau:
\choiceTF{Cỡ mẫu $n=22$}{Số trung bình là $\overline x\approx42{,}39$}{$59$ là số bao xi măng nhiều nhất bán được trong một ngày}{\True Trung vị của mẫu số liệu là $50$}
\loigiai{Sắp xếp: $21$; $33$; $36$; $36$; $36$; $43$; $43$; $45$; $45$; $45$; $47$; $50$; $50$; $50$; $53$; $54$; $54$; $58$; $59$; $61$; $62$; $65$; $67$.

a) Sai: $n=23$.

b) Sai: $\overline x=\dfrac{1113}{23}\approx48{,}39$.

c) Sai: nhiều nhất là $67$ bao.

d) Đúng: $M_e$ là giá trị thứ $12$, bằng $50$.}
\end{ex}
